\begingroup
\renewcommand{\arraystretch}{1.45}
\Needspace{23\baselineskip}
\InterpolationArticle{13}
\textsc{Theorem.}\quad\textit{If $X$ is a function of the arc $x$ of the form $(F)$}
\[
\begin{aligned}
&\alpha+\alpha'\cos x+\alpha''\cos 2x+\alpha'''\cos 3x+\text{etc.}\\[5pt]
&\phantom{\alpha}+\beta'\sin x+\beta''\sin 2x+\beta'''\sin 3x+\text{etc.}
\end{aligned}
\]
\textit{or of the form $(G)$}
\[
\begin{aligned}
&\gamma\cos\tfrac12x+\gamma'\cos\tfrac32x+\gamma''\cos\tfrac52x+\text{etc.}\\[5pt]
&+\delta\sin\tfrac12x+\delta'\sin\tfrac32x+\delta''\sin\tfrac52x+\text{etc.}
\end{aligned}
\]
\textit{and, on setting $x=a$, the value of the function $X$ becomes $=0$, then $X$ is divisible by $\sin\tfrac12(x-a)$, and the quotient is of the form $G$ in the first case and of the form $F$ in the second.}

\textit{Proof.\quad First case.}\quad If, in the function $X$, each term $\cos nx$ is replaced by $\cos nx-\cos na$, each term $\sin nx$ by $\sin nx-\sin na$,
\SourceLeaf{296}
and finally $\alpha$ by $\alpha-\alpha$ or $0$, it is evident that the constant terms added in this way cancel one another, so that the function $X$ is unchanged. Each individual term of $X$ is now divisible by $\sin\tfrac12(x-a)$, namely
\[
\begin{aligned}
\frac{\cos nx-\cos na}{\sin\tfrac12(x-a)}
&=-2\sin((n-\tfrac12)x+\tfrac12a)\\[5pt]
&\phantom{=}{-}2\sin((n-\tfrac32)x+\tfrac32a)\\[5pt]
&\phantom{=}{-}2\sin((n-\tfrac52)x+\tfrac52a)\\[5pt]
&\phantom{=}{-}\text{etc.}\\[5pt]
&\phantom{=}{-}2\sin((\tfrac12x+(n-\tfrac12)a)
\end{aligned}
\]
\[
\begin{aligned}
\frac{\sin nx-\sin na}{\sin\tfrac12(x-a)}
&=+2\cos((n-\tfrac12)x+\tfrac12a)\\[5pt]
&\phantom{=}{+}2\cos((n-\tfrac32)x+\tfrac32a)\\[5pt]
&\phantom{=}{+}2\cos((n-\tfrac52)x-\tfrac52a)\\[5pt]
&\phantom{=}{+}\text{etc.}\\[5pt]
&\phantom{=}{+}2\cos(\tfrac12x+(n-\tfrac12)a)
\end{aligned}
\]
Each coefficient can clearly be brought into the form $G$; hence the whole of $X$ is also divisible by $\sin\tfrac12(x-a)$, and the quotient can be brought into the form $G$. Q.~E.~D.

\Needspace{40\baselineskip}
\textit{Second case.}\quad If, in the function $X$, each term $\cos nx$ (where $n$ now denotes half an odd integer, rather than an integer as before) is replaced by $\cos nx-\cos\tfrac12(x-a)\cos na$, and each term $\sin nx$ by $\sin nx-\cos\tfrac12(x-a)\sin na$, it is evident that the terms added to the function $X$ in this way cancel one another, so that $X$ is unchanged. Each individual term of $X$ is now divisible by $\sin\tfrac12(x-a)$, namely
\[
\begin{aligned}
\frac{\cos nx-\cos\tfrac12(x-a)\cos na}{\sin\tfrac12(x-a)}\qquad\\[5pt]&=-2\sin((n-\tfrac12)x+\tfrac12a)\\[5pt]
&\phantom{=}{-}2\sin((n-\tfrac32)x+\tfrac32a)\\[5pt]
&\phantom{=}{-}2\sin((n-\tfrac52)x+\tfrac52a)\\[5pt]
&\phantom{=}{-}\text{etc.}\\[5pt]
&\phantom{=}{-}2\sin(x+(n-1)a)\\[5pt]
&\phantom{=}{-}\sin na
\end{aligned}
\]
\[
\begin{aligned}
\frac{\sin nx-\cos\tfrac12(x-a)\sin na}{\sin\tfrac12(x-a)}\qquad\\[5pt]&=+2\cos((n-\tfrac12)x+\tfrac12a)\\[5pt]
&\phantom{=}{+}2\cos((n-\tfrac32)x+\tfrac32a)\\[5pt]
&\phantom{=}{+}2\cos((n-\tfrac52)x+\tfrac52a)\\[5pt]
&\phantom{=}{+}\text{etc.}\\[5pt]
&\phantom{=}{+}2\cos(x+(n-1)a)\\[5pt]
&\phantom{=}{+}\cos na
\end{aligned}
\]
\SourceLeaf{297}
as is readily confirmed by multiplication. Each quotient is of the form $F$. Hence the whole function $X$ is divisible by $\sin\tfrac12(x-a)$, and the quotient can be brought into the form $F$. Q.~E.~D.

\Needspace{32\baselineskip}
\InterpolationArticle{14}
\textsc{Theorem.}\quad\textit{If a function $X$ of the form $F$ or $G$ becomes, for several distinct\HistoricalNote{24} values of $x$, namely $a$, $b$, $c$, $d$ etc., $=0$, then it is divisible by the product}
\[
\sin\tfrac12(x-a)\sin\tfrac12(x-b)\sin\tfrac12(x-c)\sin\tfrac12(x-d)\ .\,.\,.\,.
\]
\textit{and the quotient will be of the same form as $X$ or of the other form, according as the number of values $a$, $b$, $c$, $d$ $.\,.\,.\,.$ is even or odd.}

\textit{Proof.}\quad Set $X=X'\sin\tfrac12(x-a)$, and let $M$ be the value of $X'$ for $x=b$. Then $M\sin\tfrac12(b-a)$ is the value of $X$ for $x=b$; since this must be $=0$, necessarily $M=0$. Thus $X'$ is divisible by $\sin\tfrac12(x-b)$, and by the same reasoning the resulting quotient is divisible by $\sin\tfrac12(x-c)$, and so on. Hence $X$ is divisible by the product
\[
\sin\tfrac12(x-a)\sin\tfrac12(x-b)\sin\tfrac12(x-c)\sin\tfrac12(x-d)\ .\,.\,.\quad\text{Q.~E.~D.}
\]
The second part of the theorem is so evident that no proof is needed.

The converse theorem is equally evident: if $X$ is divisible by the product
\[
\sin\tfrac12(x-a)\sin\tfrac12(x-b)\sin\tfrac12(x-c)\sin\tfrac12(x-d)\ .\,.\,.\,.
\]
its value becomes $=0$ when $x$ is assigned any one of the values $a$, $b$, $c$, $d$ etc.

\endgroup
